\subsection{Primary execution settings and detailed contrasts}
The preregistered held-out study used two NVIDIA A100
80GB PCIe GPUs at a university computing facility, each running an independent
Qwen2.5-7B-Instruct instance with tensor parallelism (TP) of one. The serving stack
was vLLM 0.9.0 \citep{kwon2023vllm} with PyTorch 2.7.0+cu126, Transformers 4.51.1, tokenizers 0.21.1,
and NumPy 1.26.4. Settings were FP16, maximum model length 4096, batched-token
limit 16,384, sequence limit 1,024, four local request slots, chunked prefill,
no prefix caching, eager execution, and GPU memory utilization 0.8. Readiness
required per-process health probes and a dual-ready barrier; cleanup verified
both processes' release. Per-run environment and GPU records identify these settings.

The three held-out windows follow the earlier evaluation windows with a
30-minute exclusion interval. A predefined filter on 60 s arrival counts
identifies steady, recovery, and burst-offline classes (even, declining, and
bursty online arrivals). Deterministic hash ordering selects the first eligible
window in steady--recovery--burst order, subject to fixed separation criteria.
Online rows are thinned at a fixed threshold; each window receives 120 offline
jobs at 15 s intervals with 256-input/256-output token budgets. Selection
parameters and generated inputs are frozen before replay; the supplement
identifies the reproducibility records.


Replacing the service lookup with Vidur's profiled model tests the same
validation contract with an external predictor. Twelve retrospective CPU
rollouts retain Local's lifecycle parameters and the original guard, and reuse
all 36 physical runs. They reproduce H's recovery misses (116/120 offline
jobs on time versus 113--116 physically) and the guard's 120/120.
For H and H+guard in steady, median occupied-resource absolute percentage error (APE) is
2.47\%, yet target agreement is 33.0--35.5\%; guard recovery has 14.74\%
occupied-resource APE. Thus the external backend preserves the key deadline
contrast while revealing residual timing errors (Supplement Section~A.3).


The 14B study tests remeasured lifecycle timing in Qwen2.5-14B-Instruct/vLLM.
H fails all four runs; H+guard passes all four. Replacing 14B-Hist's historical
startup with the measured value in 14B-Local raises agreement from 4/8 to 8/8
and lowers cost MAE from 0.4105 to 0.0668 scenario USD. A retrospective
startup-only replay (20--90\,s at 1\,s, with 0.25\,s near the measured value)
yields 8/8 only at 39.25--43 and 44.25--48\,s; 43.25--44 returns 6/8 and
48.25--90 returns 4/8, all errors being false accepts of H. The fixed 30\,s
point is 6/8; 19/20 calibration samples lie in the upper tested region, while
52.35\,s gives 4/8. The finite grid is non-monotonic; event replay links
longer startup to cooldown pressure, target changes, and later deadline outcomes.

